% УВОД, без номер.
\section*{УВОД}
\label{sec:introduction}
\addcontentsline{toc}{section}{УВОД}

Сигурността на един мрежов обмен се крепи на два механизма, които потребителят не вижда:
протоколът за защита на транспортния слой и инфраструктурата на публичните ключове. И двата се
реализират от малък брой библиотеки и точно затова свойствата им рядко се проверяват
самостоятелно. Приема се, че щом връзката е установена, тя е защитена, и че щом библиотеката не
е върнала грешка, веригата от сертификати е валидна.

Проектът изхожда от обратната постановка. Договорените параметри на една връзка са наблюдаема
величина. Валидността на една верига е резултат от алгоритъм, записан в \cite{rfc5280}, който
може да бъде реализиран самостоятелно. Устойчивостта срещу лавинен трафик е свойство, което се
измерва, а не се декларира.

\textbf{Целта} е работеща и измерима система, наречена Sentinel, която чете ръкостискането по
TLS 1.3, проверява сертификатни вериги по алгоритъма от стандарта и управлява допускането на
нови връзки при претоварване, като цената на всеки механизъм е измерена, а не предположена.
\textbf{Задачите} са пет: преглед на приложимите стандарти и на заплахите; обосновка на избора
между възможните решения; проектиране на слоеста архитектура; реализация на C++20, в която
непроведените проверки се отчитат изрично като непроведени; и измерване в собствена изолирана
среда. Към тях стои изискване, което се оказа не по-малко определящо: сглобяване на чужда машина
без нито една задължителна външна зависимост.

\textbf{Обхватът} е защитен. Системата анализира, проверява и допуска. Тя не съдържа средства за
проникване, а опитите с лавинен трафик се провеждат срещу собствен стенд, чийто слушател се
свързва само с адреса на локалния интерфейс.

\textbf{Ограничение, което важи за целия документ.} Системата не е подлагана на атака, на фъзинг
или на външен одит. Всяко твърдение по-нататък се отнася до това, което кодът проверява, отчита
и измерва, а не до доказана устойчивост срещу нападател.
